\subsection{Modular Arithmetic}\label{sec:mathematical_underpinnings:modular_arithmetic}

Modular arithmetic is a system for arithmetic operations where numbers (or expressions) \textit{wrap around} upon reaching a certain fixed number (or expression) called \textbf{modulus}. For instance, numbers are reduced by the modulus repeatedly as to reach a number between 0 and the modulus itself. A generic value $x$ under a modulus $n$ is usually written as ``$x \pmod{n}$'', and it is equal to the remainder of $x$ upon division by $n$. An operation $x \bigcdot y$ executed with a modulus $n$ is usually written $(x \bigcdot y) \% n$ or, more generally, $x \bigcdot y \pmod{n}$.

\exampleBlock{An example of real-world modular arithmetic}{On a 12-hour clock, if it's 9 o'clock and we add 5 hours, we get $9 + 5 = 14 \equiv 2 \pmod{12}$.\footnote{While both symbols represent relationships between expressions, the equal sign ($=$) generally represents a direct value-for-value equality, whereas the triple bar ($\equiv$) means that two expressions are equivalent or congruent in all aspects --- not just in value --- in a given context (in this case, $14 \equiv 2$ in the context of modulus $12$).}}

Modular arithmetic is useful to both limit sets which otherwise would be infinite. For instance, the set of integers $\mathbb{Z} = \{ \ldots, -2, -1, 0, 1, 2, \ldots \}$ is infinite, while the set of integers modulo $12$ is $\mathbb{Z}_{12} = \{ 0, 1, \ldots , 10, 11 \}$ and counts exactly 12 elements.

\curiosityBlock{The elements of $\mathbb{Z}_{12}$ are not actually integers}{More formally, the 12 elements in $\mathbb{Z}_{12}$ are so-called \textbf{equivalence classes} of integers modulo $12$ --- not single integers. Informally, the equivalence class of an integer $x$ modulo $12$ is a set containing all integers that leave the same remainder upon division by 12; for instance, the integers $2, 14, 26, 38$, and $50$ are part of the same equivalence class (as they all have remainder 2 upon division by 12). A common notation for representing equivalence classes are square parentheses (i.e.,$\mathbb{Z}_{12} = \{ [0], [1], \ldots , [10], [11] \}$), where $[x] = \{ z \in \mathbb{Z} \sth z \equiv x \pmod{12} \}$. Another notation you may see in this context is a horizontal bar on top of an integer; for instance, $\overline{x}$ denotes the residue class of $x$ modulo $12$ (a \textbf{residue class} is the set of all integers that leave the same remainder upon division by a modulus). Hence, $[x] \equiv \overline{x}$ in modular arithmetic. However, note that bars may denote other concepts in different contexts, so their use may sometimes be ambiguous. Summarizing, the set of integers modulo $12$ should be written as $\mathbb{Z}_{12} = \{ [0], [1], \ldots , [10], [11] \}$ or as $\mathbb{Z}_{12} = \{ \overline{0}, \overline{1}, \ldots , \overline{10}, \overline{11} \}$. Nonetheless, for simplicity, it is common to informally just write $\mathbb{Z}_{12} = \{ 0, 1, \ldots , 10, 11 \}$.}

Notably, using modular arithmetic to limit sets makes reversing some operations difficult, e.g., the \gls{dl} problem (\Cref{sec:trapdoor_functions:discrete_logarithm}) --- the underlying intuition is that, every time the result of an operation exceeds the modulus and therefore \textit{wraps around} it, some information is lost. For instance, $2 \pmod{12}$ may be the result of both $9 + 5 \pmod{12}$ but also $7 + 7 \pmod{12}$; while computing an addition modulo $12$ is easy, reversing it --- i.e., start from $2$ and then go back to the addends --- is difficult because more choices may be available.
